% Source-derived chapter 01: Introduction
% Original source: squarefree-values-polynomial-discriminants
\chapter{Introduction}
\label{squarefree-values-polynomial-discriminants-chapter-01}
\source{squarefree-values-polynomial-discriminants}

\begin{remark}[Source section]
\label{squarefree-values-polynomial-discriminants-chapter-01-source}
\source{squarefree-values-polynomial-discriminants}
\source{squarefree-values-polynomial-discriminants}
This chapter reproduces the corresponding section of the retrieved
LaTeX source, including its definitions, statements, examples, and
proofs. It has not yet been translated into Lean.
\notready
\end{remark}

% The source section title was: Introduction
The pupose of this paper is to determine the density of monic integer
polynomials of given degree whose discriminant is squarefree.  For
polynomials $f(x)=x^n+a_1x^{n-1}+\cdots+a_n$, the term $(-1)^ia_i$
represents the sum of the $i$-fold products of the roots of $f$. It is
thus natural to order monic polynomials
$f(x)=x^n+a_1x^{n-1}+\cdots+a_n$ by the height
$H(f):=\max\{|a_i|^{1/i}\}$ (see, e.g., \cite{BG2}, \cite{PS2},
\cite{SW}).  We determine the density of monic integer polynomials
of degree $n$ having squarefree discriminant with respect to the ordering by this
height, and show that the density is positive.  The existence of infinitely many monic integer
polynomials of each degree having squarefree discriminant was
first demonstrated by Kedlaya~\cite{Kedlaya}.
However, it has not previously been known whether the density exists or even that the lower density
is positive.

To state the theorem, define the constants $\lambda_n(p)$ by
\begin{equation}\label{jos}
\source{squarefree-values-polynomial-discriminants}
\lambda_n(p)=\left\{
\begin{array}{cl}
1 & \mbox{if $n =1$,}\\[.0875in]
1-\displaystyle\frac1{p^2} & \mbox{if $n= 2$,}\\[.1875in]
%1-\displaystyle\frac2{p^2}+\frac1{p^3} & \mbox{if $n= 3$,}\\[.185in]
%1-\displaystyle\frac1{p}+\frac{(p-1)^2(1-(-p)^{2-n})}{p^2(p+1)} & \mbox{if $n\geq 4$}\\
1-\displaystyle \frac{3p^{n-1}-p^{n-2}+(-1)^n(p-1)^2}{p^n(p+1)} & \mbox{if $n\geq 3$}
\end{array}\right.
\end{equation}
for $p\neq 2$; also, let $\lambda_1(2)=1$ and $\lambda_n(2)=1/2$ for
$n\geq2$.  Then a result of Yamamura~\cite[Proposition~3]{Yamamura2}
%(see also Brakenhoff~\cite[Theorem~6.9]{ABZ})
states
that $\lambda_n(p)$ is the density of monic polynomials of degree $n$ over~$\Z_p$
having discriminant indivisible by~$p^2$.
Let~$\lambda_n:=\prod_p\lambda_n(p)$, where the product is over all
primes $p$.  We~prove:

\begin{theorem}
\source{squarefree-values-polynomial-discriminants}
\notready\label{polydisc2}
\source{squarefree-values-polynomial-discriminants}
Let $n\geq1$ be an integer.  Then when
monic integer polynomials $f(x)=x^n+a_1x^{n-1}+\cdots+a_n$ of
degree~$n$ are ordered by $H(f):=
\max\{|a_1|,|a_2|^{1/2},\ldots,|a_n|^{1/n}\}$, the density having
squarefree discriminant $\Delta(f)$ exists and is equal to $\lambda_n>0$.
\end{theorem}
Our method of proof implies that the theorem remains true even if we
restrict only to those polynomials of a given degree $n$ having a
given number of real roots.

It is easy to see from the definition of the $\lambda_n(p)$ that the
$\lambda_n$ rapidly approach a limit $\lambda$ as $n\to\infty$,
namely,
\begin{equation}
\lambda=\lim_{n\to\infty} \lambda_n = \frac12 \cdot \prod_{p\geq 3}
%\left(1-\displaystyle\frac1{p}+\frac{(p-1)^2}{p^2(p+1)}\right)
\left(1-\displaystyle\frac{3p-1}{p^2(p+1)}\right)
\approx 30.7056\%.
\end{equation}
Therefore, as the degree tends to infinity, the
probability that a random monic integer polynomial has squarefree
discriminant tends to $\lambda\approx 30.7056\%$.

In algebraic number theory, one often considers number fields that are
defined as a quotient ring $K_f:=\Q[x]/(f(x))$ for some irreducible
integer polynomial $f(x)$. The question naturally arises as to whether
$R_f:=\Z[x]/(f(x))$ gives the ring of integers of $K_f$.  Our second
main theorem states that this is in fact the case for {most}
polynomials $f(x)$. We prove:

\begin{theorem}
\source{squarefree-values-polynomial-discriminants}
\notready\label{polydiscmax2}
\source{squarefree-values-polynomial-discriminants}
  Let $n > 1$ be an integer.  Then when monic integer polynomials
  $f(x)=x^n+a_1x^{n-1}+\cdots+a_n$ of degree~$n$ are ordered by
  $H(f)$, the density of polynomials $f$ such that $\Z[x]/(f(x))$ is
  the ring of integers in its fraction field is
  $\prod_p(1-1/p^2)=\zeta(2)^{-1}$.
\end{theorem}
%\begin{theorem}\label{polydiscmax2}
%Let $n\geq1$ be an integer.  The density of irreducible monic integer
%polynomials $f(x)=x^n+a_1x^{n-1}+\cdots+a_n$ of degree~$n>1$, when
%ordered by $H(f)$, such that $\Z[x]/(f(x))$ is the ring of integers in
%its fraction field is $\prod_p(1-1/p^2)=\zeta(2)^{-1}$.
%\end{theorem}
Note that $\zeta(2)^{-1}\approx\, 60.7927\%$.  Since a density of
100\% of monic integer polynomials are irreducible (and indeed have
associated Galois group $S_n$) by Hilbert's irreducibility theorem, it
follows that $\approx 60.7927\%$ of monic integer polynomials $f$ of
any given degree $n>1$ have the property that $f$ is irreducible and
$\Z[x]/(f(x))$ is the maximal order in its fraction field. The
quantity
\begin{equation}\label{eqrho}
\source{squarefree-values-polynomial-discriminants}
  \rho_n(p):=1-\frac{1}{p^2}
\end{equation}
represents the density of monic integer polynomials
of degree $n>1$ over $\Z_p$ such that $\Z_p[x]/(f(x))$ is the maximal
order in $\Q_p[x]/(f(x))$.  The determination of this beautiful
$p$-adic density, and its independence of $n$, is due to Hendrik
Lenstra (see~\cite[Proposition~3.5]{ABZ}).  Theorem~\ref{polydiscmax2}
again holds even if we restrict to polynomials of degree $n$ having a
fixed number of real roots.

If the discriminant of an order in a number field is squarefree, then
that order must be maximal.  Thus the irreducible polynomials counted
in Theorem~\ref{polydisc2} are a subset of those counted in
Theorem~\ref{polydiscmax2}.  The additional usefulness of
Theorem~\ref{polydisc2} in some arithmetic applications is that if
$f(x)$ is a monic irreducible integer polynomial of degree $n$ with
squarefree discriminant, then not only is $\Z[x]/(f(x))$ maximal in
the number field $\Q[x]/(f(x))$ but the associated Galois group is
necessarily the symmetric group $S_n$ (see, e.g., \cite{Scholz},
\cite{Yamamura}, \cite{Nakagawa}, \cite{Kondo} for further details and
applications).

We prove both Theorems~\ref{polydisc2} and \ref{polydiscmax2} with
power-saving error terms. More precisely, let $\Vmn(\Z)$ denote
the subset of $\Z[x]$ consisting of all monic integer polynomials of
degree $n$.  Then it is easy to see that
\begin{equation*}
{\#\{f\in \Vmn(\Z):  H(f)<X \}} = 2^nX^{\frac{\scriptstyle n(n+1)}{\scriptstyle2}} + O(X^{\frac{\scriptstyle n(n+1)}{\scriptstyle2}-{ 1}}).
\end{equation*}
We prove
\begin{equation}\label{errorterms}
\source{squarefree-values-polynomial-discriminants}
{\begin{array}{ccl}
\displaystyle
{\#\{f\in \Vmn(\Z) : H(f)<X \mbox{ and $\Delta(f)$ squarefree}\}}
&\!\!=\!\!&  \lambda_n\cdot 2^n{X^\frac{\scriptstyle n(n+1)}{\scriptstyle2}} + O_\varepsilon(X^{\frac{\scriptstyle n(n+1)}{\scriptstyle2}-{\textstyle \frac15}+\varepsilon});\\[.15in]
\displaystyle
{\#\{f\in \Vmn(\Z) : H(f)<X \mbox{ and $\Z[x]/(f(x))$ maximal}\}}&\!\!=\!\!&  {\displaystyle\frac{6}{\pi^2}}\cdot 2^nX^{\frac{\scriptstyle n(n+1)}{\scriptstyle2}} + O_\varepsilon(X^{\frac{\scriptstyle n(n+1)}{\scriptstyle2}-{\textstyle \frac15}+\varepsilon})
\end{array}}
\end{equation}
for $n>1$.

These asymptotics imply Theorems~\ref{polydisc2} and
\ref{polydiscmax2}.  Since it is known that the number of reducible
monic polynomials of a given degree~$n$ is of a strictly smaller order
of magnitude than the error terms above (see Proposition
\ref{propredboundall}), it does not matter whether we require $f$ to
be irreducible in the above asymptotic formulae.

Recall that a number field $K$ is called {\it monogenic} if its ring
of integers is generated over~$\Z$ by one element, i.e., if
$\Z[\theta]$ gives the maximal order of $K$ for some $\theta\in K$.
As a further application of our methods, we obtain the following
corollary to Theorem~\ref{polydisc2}:

\begin{corollary}
\source{squarefree-values-polynomial-discriminants}
\notready\label{monogenic}
\source{squarefree-values-polynomial-discriminants}
Let $n>1$.  The number of isomorphism classes of number fields
of degree~$n$ and absolute discriminant less than $X$ that are
monogenic and have associated Galois group $S_n$ is $\gg X^{1/2+1/n}$.
\end{corollary}
We note that our lower bound for the number of monogenic $S_n$-number
fields of degree $n$ improves slightly the best-known lower bounds for
the number of $S_n$-number fields of degree $n$, due to Ellenberg and
Venkatesh~\cite[Theorem 1.1]{EV}, by simply forgetting the
monogenicity condition in Corollary~\ref{monogenic}.  We conjecture
that the exponent in our lower bound in Corollary~\ref{monogenic} for
monogenic number fields of degree $n$ is optimal.

As is illustrated by Corollary~\ref{monogenic},
Theorems~\ref{polydisc2} and \ref{polydiscmax2} give a powerful method
to produce number fields of a given degree having given properties or
invariants.  We give one further example of interest.  Given a number
field $K$ of degree $n$ with $r$ real embeddings $\xi_1,\dots,\xi_r$
and $s$ complex conjugate pairs of complex embeddings
$\xi_{r+1},\bar\xi_{r+1},\ldots,\xi_{r+s},\bar\xi_{r+s}$, the ring of
integers $\mathcal O_K$ may naturally be viewed as a lattice in $\R^n$
via the map $x\mapsto (\xi_1(x),\ldots,\xi_{r+s}(x))\in
\R^r\times\C^s\cong \R^n$.  We may thus ask about the length of the
shortest vector in this lattice generating $K$.

In their final remark~\cite[Remark~3.3]{EV}, Ellenberg and Venkatesh
conjecture that the number of number fields $K$ of degree $n$ whose
shortest vector in $\OO_K$ generating $K$ is of length less than~$Y$
is $\,\asymp Y^{(n-1)(n+2)/2}$. They prove an upper bound of this
order of magnitude.  We use Theorem~\ref{polydiscmax2} to prove also a
lower bound of this size, thereby proving their conjecture:

\begin{corollary}
\source{squarefree-values-polynomial-discriminants}
\notready\label{shortvector}
\source{squarefree-values-polynomial-discriminants}
Let $n>1$.  The number of isomorphism classes of number fields $K$ of
degree~$n$ whose shortest vector in
%the ring of integers
$\OO_K$ generating $K$ has length less than $Y$ is $\,\asymp$
$Y^{(n-1)(n+2)/2}$.  The same is true if we further impose the
condition that the Galois group of the normal closure of $K$ is $S_n$.
\end{corollary}
%\begin{corollary}\label{shortvector}
%Let $n>1$.  The number of
%isomorphism classes of number fields $K$ of degree~$n$ whose shortest
%vector in $\OO_K$ generating $K$ has length less than $Y$ is
%$\,\asymp$ $Y^{(n-1)(n+2)/2}$.
%\end{corollary}

%Again,
%Corollary~\ref{shortvector} remains true even if we impose the
%condition that the associated Galois group is~$S_n$ (by using
%Theorem~\ref{polydisc2} instead of Theorem~\ref{polydiscmax2}).

Finally, we remark that our methods allow the analogues of all of the
above results to be proven with any finite set of local conditions
imposed at finitely many places (including at infinity); the orders of
magnitudes in these theorems are then seen to remain the same---with
different (but easily computable in the cases of
Theorems~\ref{polydisc2} and \ref{polydiscmax2}) positive
constants---provided that no local conditions are imposed that force
the set being counted to be empty (i.e., no local conditions are
imposed at~$p$ in Theorem~\ref{polydisc2} that force $p^2$ to divide
the discriminant, no local conditions are imposed at~$p$ in
Theorem~\ref{polydiscmax2} that cause $\Z_p[x]/(f(x))$ to be
non-maximal over~$\Z_p$, and no local conditions are imposed at $p$ in
Corollary~\ref{monogenic} that cause such number fields to be
non-monogenic locally). In fact, we can even impose certain infinite
sets of local conditions (see Theorem \ref{thmaingencong}).



\vspace{.1in} We now briefly describe our methods.  It is easily seen
that the desired densities in Theorems~\ref{polydisc2} and
\ref{polydiscmax2}, if they exist, must be bounded above by the Euler
products $\prod_p \lambda_n(p)$ and $\prod_p (1-1/p^2)$,
respectively. The difficulty is to show that these Euler products are
also the correct lower bounds.  As is standard in sieve theory, to
demonstrate the lower bound, a ``tail estimate'' is required to show
that not too many discriminants of polynomials $f$ are divisible by
$p^2$ when $p$ is large relative to the discriminant $\Delta(f)$ of
$f$ (here, large means larger than $\Delta(f)^{1/(n-1)}$, say).

For any prime $p$, and a monic integer polynomial $f$ of degree $n$
such that $p^2\mid \Delta(f)$, we say that $p^2$ {\it strongly
  divides} $\Delta(f)$ if $p^2\mid \Delta(f + pg)$ for any integer
polynomial $g$ of degree $n$; otherwise, we say that $p^2$ {\it weakly
  divides} $\Delta(f)$. Then $p^2$ strongly divides $\Delta(f)$ if and
only if $f$ modulo $p$ has at least two distinct multiple roots in
$\bar{\F}_p$, or has a root in $\F_p$ of multiplicity at least 3; and
$p^2$ weakly divides $\Delta(f)$ if $p^2\mid \Delta(f)$ but $f$ modulo
$p$ has only one multiple root in $\F_p$ and this root is a simple
double root.

For any squarefree positive integer $m$, let $\U_m^\s$
(resp.\ $\U_m^\w$) denote the set of monic integer polynomials in
$\Vmn(\Z)$ whose discriminant is strongly divisible (resp.\ weakly
divisible) by $p^2$ for every prime factor $p$ of $m$.  Then we prove
tail estimates for $\U_m^\s$ and $\U_m^\w$ separately, as follows.

\begin{theorem}
\source{squarefree-values-polynomial-discriminants}
\notready\label{thm:mainestimate}
\source{squarefree-values-polynomial-discriminants}
For any positive real number $M$ and any $\epsilon>0$, we have

\vspace{-5pt}\begin{eqnarray*}
\label{eq:equs}
\source{squarefree-values-polynomial-discriminants}
{\rm (a)}\quad \#\bigcup_{\substack{m>M\\ m\;\mathrm{ squarefree}
 }}\{f\in\U_m^\s:H(f)<X\}&=&
O_\epsilon(X^{n(n+1)/2+\epsilon}/M)+O(X^{n(n+1)/2-1});\\[.075in]
\label{equ1}
\source{squarefree-values-polynomial-discriminants}
{\rm (b)}\quad
\#\bigcup_{\substack{m>M\\
m\;\mathrm{ squarefree}
}}\{f\in\U_m^\w:H(f)<X\}&=&
O_\epsilon(X^{n(n+1)/2+\epsilon}/M)+O_\epsilon(X^{n(n+1)/2-1/5+\epsilon}),
\end{eqnarray*}
where the implied constants are independent of $M$ and $X$.
\end{theorem}
\noindent To prove our main theorems, we will use Theorem
\ref{thm:mainestimate} with $M=X^{1/2}$.


The power savings in the error terms above also have applications towards
determining the distributions of low-lying zeros in families of
Dedekind zeta functions of monogenic degree-$n$ fields; see \cite[\S5.2]{SST}.

We prove the estimate in the strongly divisible case (a) of
Theorem~\ref{thm:mainestimate} by geometric techniques, namely, a
quantitative version of the Ekedahl sieve (\cite{Ek}, \cite[Theorem
  3.3]{geosieve}). While the proof of \cite[Theorem~3.3]{geosieve}
uses homogeneous heights, and considers the union over all primes
$p>M$, the same proof also applies in our case of weighted homogeneous
heights, and a union over all squarefree $m>M$. Since the last
coefficient $a_n$ is in a larger range than the other coefficients, we
in fact obtain a smaller error term than in \cite[Theorem~3.3]{geosieve}.

The estimate in the weakly divisible case (b) of
Theorem~\ref{thm:mainestimate} is considerably more difficult. Our
main idea is to embed polynomials $f$, whose discriminant is {\it weakly}
divisible by $p^2$, into a larger space that has more symmetry, such
that the invariants under this symmetry are given exactly by the
coefficients of $f$; moreover, we arrange for the image of $f$ in the
bigger space to have discriminant {\it strongly} divisible by $p^2$.  We
then count in the bigger space.

More precisely, we make use of the representation of $G=\SO_n$ on the
space $W=W_n$ of symmetric $n\times n$ matrices, as studied in
\cite{BG2,SW}. We fix $A_0$ to be the $n\times n$ symmetric matrix
with $1$'s on the anti-diagonal and $0$'s elsewhere. The group
$G=\SO(A_0)$ acts on $W$ via the action $g\cdot B=gBg^t$ for $g\in G$
and $B\in W$.  Define the {\it invariant polynomial} of an element
$B\in W$ by $$f_B(x) = (-1)^{n(n-1)/2}\det(A_0x - B).$$ Then $f_B$ is
a monic polynomial of degree~$n$. It is known (see \cite[\S 4]{AITBG})
that the ring of polynomial invariants for the action of $G$ on $W$ is
freely generated by the coefficients of the invariant
polynomial. Define the {\it discriminant} $\Delta(B)$ and {\it height}
$H(B)$ of an element $B\in W$ by $\Delta(B)=\Delta(f_B)$ and
$H(B)=H(f_B)$.  This representation of $G$ on $W$ was used in
\cite{BG2,SW} to study 2-descent on the hyperelliptic curves
$C:y^2=f_B(x)$.

A key step of our proof of Theorem~\ref{thm:mainestimate}(b) is the
construction, for every positive squarefree integer $m$, of a map
\begin{equation*}
\sigma_m:\U_m^\w\to \frac14W(\Z),
\end{equation*}
such that $f_{\sigma_m(f)}=f$ for every $f\in \U_m^\w$; here
$\frac14W(\Z)\subset W(\Q)$ is the lattice of elements $B$ whose
coefficients have denominators dividing $4$.  In our construction, the
image of $\sigma_m$ in fact lies in a special subspace $W_0$ of $W$;
namely, if $n=2g+1$ is odd, then $W_0$ consists of symmetric matrices
$B\in W$ whose top left $g\times g$ block is 0, and if $n=2g+2$ is
even, then $W_0$ consists of symmetric matrices $B\in W$ whose top
left $g\times (g+1)$ block is 0. We associate to any element of $W_0$
a further polynomial invariant which we call the $Q$-{\it invariant}
(which is a relative invariant for the subgroup of $\SO(A_0)$ that
fixes $W_0$).
%The significance of the $Q$-invariant is that, if the
%discriminant polynomial $\Delta$ is restricted to $W_0$, then it is
%not irreducible as a polynomial in the coordinates of $W_0$, but rather is
%divisible by the polynomial $Q^2$.
Next, we show that for elements $B$ in the image of $\sigma_m$, we
have $|Q(B)|=m$.
Finally, {even though the discriminant polynomial of $f\in\U_m^\w$ is
  {\it weakly} divisible by $p^2$, the discriminant polynomial of its
  image $\sigma_m(f)$, when viewed as a polynomial on $W_0\cap
  \frac14W(\Z)$, is {\it strongly} divisible by $p^2$.}  This is the
key point of our construction.

To obtain Theorem~\ref{thm:mainestimate}(b), it thus suffices to
estimate the number of $G(\Z)$-equivalence classes of elements $B\in
W_0\cap \frac14W(\Z)$ of height less than $X$ having $Q$-invariant
larger than $M$.  This can be reduced to a geometry-of-numbers
argument in the spirit of \cite{BG2,SW}, although the current count is
more subtle in that we are counting certain elements in a cuspidal
region of a fundamental domain for the action of $G(\Z)$ on $W(\R)$.
The $G(\Q)$-orbits of elements $B\in W_0\cap W(\Q)$ are called {\it
  distinguished orbits} in \cite{BG2,SW}, as they correspond to the
identity 2-Selmer elements of the Jacobians of the corresponding
hyperelliptic curves $y^2=f_B(x)$ over $\Q$; these were not counted
separately by the geometry-of-numbers methods of \cite{BG2,SW}, as
these elements lie deeper in the cusps of the fundamental domains.  We
develop a method to count those elements in the cusp having bounded height and $Q$-invariant larger than $M$, following
the arguments of \cite{BG2,SW} while using the invariance and
algebraic properties of the $Q$-invariant polynomial.  This yields
Theorem~\ref{thm:mainestimate}(b), which then allows us to carry out
the sieves required to obtain Theorems~\ref{polydisc2} and
\ref{polydiscmax2}.

Corollary~\ref{monogenic} can be deduced from
Theorem~\ref{polydisc2} roughly as follows. Let $g\in \Vmn(\R)$ be
a monic real polynomial of degree~$n$ and nonzero discriminant having
$r$ real roots and $2s$ complex roots.  Then $\R[x]/(g(x))$ is
isomorphic to $\R^n\cong \R^r\times \C^s$ via its real and complex
embeddings.  Let $\theta$ denote the image of $x$ in $\R[x]/(g(x))$
and let $R_g$ denote the lattice formed by taking the $\Z$-span of
$1,\theta,\ldots,\theta^{n-1}$. Suppose further that there exist
monic integer polynomials $h_i$ of degree $i$ for $i=1,\ldots,n-1$
such that $1,h_1(\theta),h_2(\theta),\ldots,h_{n-1}(\theta)$ is the
unique Minkowski-reduced basis of $R_g$; we say that the polynomial
$g(x)$ is {\it strongly quasi-reduced} in this case. Note that if $g$
is an integer polynomial, then the lattice $R_g$ is simply the image
of the ring $\Z[x]/(g(x))\subset \R[x]/(g(x))$ in $\R^n$ via its
archimedean embeddings.

  When ordered by their heights, we prove that 100\% of monic integer
  polynomials $g(x)$ are strongly quasi-reduced. We furthermore prove
  that two distinct strongly quasi-reduced integer polynomials $g(x)$
  and $g^\ast(x)$ of degree~$n$ with vanishing $x^{n-1}$-term
  necessarily yield non-isomorphic rings $R_g$ and $R_{g^\ast}$.  The
  proof of the positive density result of Theorem~\ref{polydisc2} then
  produces $\gg X^{1/2+1/n}$ strongly quasi-reduced monic integer
  polynomials $g(x)$ of degree~$n$ having vanishing $x^{n-1}$-term,
  squarefree discriminant, and height less than $X^{1/(n(n-1))}$.
  These therefore correspond to $\gg X^{1/2+1/n}$ non-isomorphic
  monogenic rings of integers in $S_n$-number fields of degree $n$
  having absolute discriminant less than~$X$, and
  Corollary~\ref{monogenic} follows.

A similar argument proves Corollary~\ref{shortvector}. Suppose $f(x)$
is a strongly quasi-reduced irreducible monic integer polynomial of
degree $n$ with squarefree discriminant $\Delta(f)$.  Elementary
estimates show that if $H(f)<Y$, then $\|\theta\|\ll Y$, and so the
shortest vector in the ring of integers generating the field also has
length bounded by $O(Y)$. The above-mentioned result on the number of
strongly quasi-reduced irreducible monic integer polynomial of degree
$n$ with squarefree discriminant, vanishing $x^{n-1}$-coefficient, and
height bounded by $Y$ then gives the desired lower bound of
%Corollary~\ref{shortvector}.
$\gg Y^{(n-1)(n+2)/2}.$ We give full proofs of Corollaries \ref{monogenic} and
\ref{shortvector} in Section~5.

In a subsequent paper \cite{SqSieve2} (Part II), we prove the
corresponding results for non-monic integer polynomials of degree $n$
ordered by the maximum of the absolute values of the coefficients.
Namely, we determine the density of such polynomials 
having squarefree discriminant and the density corresponding to maximal orders. The
treatment of non-monic integer polynomials in Part II builds on the
ideas here, but involves a number of new ideas due to the fact that
there exist non-monic integer polynomials $f(x)$ of degree $n$ for
which there are no symmetric integer matrices $A$ and $B$ such that
$f(x) = \det(Ax-B)$!  (See \cite{BGWhyper}.)  This
complication requires additional methods to adapt the proof to the
non-monic case.  The non-monic case has a number of applications as
well, including new results towards counting number fields and also to
the resolution of a conjecture of Poonen \cite{PBertini} on an
arithmetic Bertini theorem for the projective line. The main
theorems of this paper are the analogous results for the affine
line. The results in \cite{PBertini}, building on work of
Granville~\cite{Gabc}, imply versions of Theorems \ref{polydisc2} and
\ref{polydiscmax2} conditional on the ABC Conjecture.

\vspace{.1in} 
This paper is organized as follows. In Sections~\ref{sec:monicodd} and~\ref{sec:moniceven}, we begin by
collecting some algebraic facts about the representation $2\otimes
g\otimes(g+1)$ of $\SL_2\times\GL_g\times\GL_{g+1}$ and we define the
$Q$-invariant, which is a relative polynomial invariant for this
action.  We
then apply geometry-of-numbers techniques as described above to prove
the critical estimates of Theorem~\ref{thm:mainestimate}, handling the cases of $n$~odd and $n$~even separately. In
Section~\ref{sec:sieve}, we then show how our main theorems, Theorems~\ref{polydisc2} and \ref{polydiscmax2}, can be deduced from
Theorem~\ref{thm:mainestimate}.  Finally, in Section~\ref{latticearg},
we prove Corollary~\ref{monogenic} on the number of monogenic
$S_n$-number fields of degree~$n$ having bounded absolute
discriminant, as well as Corollary~\ref{shortvector} on the number of
rings of integers in number fields of degree $n$ whose shortest vector
generating the number field is of bounded length.
